\section{Trade-off Analysis of LoRA}

\subsection{Overview}

LoRA (Low-Rank Adaptation) is a parameter-efficient fine-tuning
method that reduces the number of trainable parameters by freezing
the original pretrained weights and introducing trainable
low-rank matrices.

Although LoRA offers significant advantages in training efficiency
and memory usage, it also introduces trade-offs in model quality,
training time, inference efficiency, and experimental reliability.

This section analyzes these trade-offs and identifies situations
where LoRA may not be the most suitable approach.

\subsection{Model Quality}

LoRA can achieve competitive performance compared to Full
Fine-Tuning while updating only a small fraction of the original
model parameters.

However, its adaptation capacity depends on the selected rank
and the placement of adapter modules.

A low rank may limit the model's ability to represent complex
weight updates, potentially reducing performance on tasks that
require substantial changes to the pretrained model.

Increasing the rank provides greater adaptation capacity but
also increases the number of trainable parameters.

Therefore, LoRA does not guarantee the same performance as
Full Fine-Tuning in every task. Model quality must be evaluated
using appropriate validation and test metrics.

\subsection{Memory Usage}

One of LoRA's main advantages is its reduced memory requirement
during training.

By freezing the original pretrained weights, LoRA avoids
storing gradients and optimizer states for those frozen
parameters.

The calculation in RE-03 illustrates this advantage for a
$1024 \times 1024$ weight matrix:

\begin{itemize}
  \item Full Fine-Tuning: 1,048,576 trainable parameters.
  \item LoRA with rank 4: 8,192 trainable parameters.
  \item LoRA with rank 8: 16,384 trainable parameters.
  \item LoRA with rank 16: 32,768 trainable parameters.
\end{itemize}

For rank 4, the number of trainable parameters is reduced
by 99.21875\% compared to Full Fine-Tuning.

However, this reduction does not imply an equivalent reduction
in total GPU memory consumption.

The original model weights must still be stored, while
activations, temporary tensors, and computational overhead
continue to consume memory.

Actual memory usage also depends on precision, batch size,
sequence length, model architecture, and optimizer configuration.

\subsection{Training Time}

LoRA can reduce training time because fewer parameters require
gradient computation, optimizer updates, and optimizer-state
storage.

Nevertheless, the original model still participates in
forward computation, and backward computation is still
required to obtain gradients for the trainable adapters.

Consequently, training speed improvements depend on the model,
hardware, dataset, batch size, and implementation.

LoRA should not be assumed to be faster in every environment.
For small models or short training runs, the time savings
may be relatively limited.

\subsection{Inference Efficiency}

LoRA introduces additional low-rank computations when its
adapters remain separate from the original weights.

The updated weight matrix can be expressed as:

\[
W' = W + BA
\]

where $W$ is the frozen pretrained weight matrix and
$A$ and $B$ are the trainable low-rank matrices.

To reduce additional inference overhead, the adapter weights
can be merged into the original weights:

\[
W_{\text{merged}} = W + BA
\]

After merging, inference can use the updated weight matrix
without separately applying the adapter computation.

However, merging reduces the flexibility of independently
switching between adapters and may require additional memory
during the merging process.

Therefore, inference efficiency depends on the deployment
strategy and whether adapters are merged.

\subsection{Experimental Risks and Limitations}

Several factors can affect the reliability and reproducibility
of LoRA experiments.

\begin{itemize}
  \item \textbf{Hyperparameter sensitivity:}
  Results depend on rank, learning rate, batch size, and
  other training configurations.

  \item \textbf{Rank selection:}
  An excessively low rank may limit adaptation capacity,
  while a high rank increases computational and memory costs.

  \item \textbf{Dataset quality:}
  Limited, noisy, or imbalanced training data may lead
  to poor generalization.

  \item \textbf{Overfitting:}
  The model may perform well on training data but fail
  to generalize to unseen examples.

  \item \textbf{Reproducibility:}
  Different random seeds, hardware, software versions,
  and configurations may produce different results.

  \item \textbf{Evaluation limitations:}
  A single experiment or evaluation metric is insufficient
  to establish that LoRA is consistently better than
  Full Fine-Tuning.
\end{itemize}

A reliable comparison should use the same datasets,
evaluation metrics, and comparable experimental settings.
Recording training time and actual GPU memory consumption
is also important.

Where possible, multiple random seeds should be used
to evaluate result stability.

\subsection{When LoRA May Not Be the Best Choice}

Despite its advantages, LoRA is not always the optimal
fine-tuning method.

The following situations may favor Full Fine-Tuning
or alternative adaptation methods:

\begin{itemize}
  \item \textbf{Extensive model adaptation:}
  Tasks requiring substantial changes across many model
  parameters may benefit from Full Fine-Tuning.

  \item \textbf{Insufficient adaptation capacity:}
  If LoRA with practical rank settings cannot achieve
  the required performance, Full Fine-Tuning or higher-rank
  adaptation may be preferable.

  \item \textbf{Small models with sufficient resources:}
  When GPU memory and computational resources are not
  significant constraints, Full Fine-Tuning may offer
  a simpler training workflow.

  \item \textbf{Strict inference requirements:}
  Applications with strict latency or deployment constraints
  may prefer merged adapters or directly fine-tuned models.

  \item \textbf{Unsatisfactory experimental results:}
  If LoRA consistently fails to meet the required quality
  after reasonable hyperparameter tuning, alternative
  fine-tuning strategies should be considered.
\end{itemize}

These situations do not mean that LoRA always performs
worse. The appropriate method depends on the task,
available resources, and deployment requirements.

\subsection{Overall Discussion}

The results from RE-03 demonstrate that LoRA can
substantially reduce the number of trainable parameters.

The experimental findings from RE-04 and RE-05 should
be considered alongside these calculations to evaluate
model quality, training efficiency, and practical limitations.

In particular, parameter reduction alone is not sufficient
to determine the overall effectiveness of LoRA.

A complete evaluation must consider model performance,
actual GPU memory consumption, training time, inference
efficiency, and experimental reliability.

\subsection{Conclusion}

LoRA provides an effective parameter-efficient alternative
to Full Fine-Tuning, particularly when computational
resources and GPU memory are limited.

Its main advantages are reduced trainable parameters,
lower gradient and optimizer-state requirements, and
the flexibility of using separate adapters.

However, these benefits involve trade-offs in adaptation
capacity, training performance, and deployment complexity.

LoRA is therefore not universally superior to Full
Fine-Tuning. The final choice should be based on
experimental evidence and the specific requirements
of the application.
